\subsection{Recursive Best First Search (GBFS)}

\textbf{Ý tưởng:}
\begin{itemize}
    \item \textbf{Recursive Best-First Search (RBFS)} là một phiên bản \textbf{tối ưu về bộ nhớ} của thuật toán A* \cite{lecture2025}.  
    Nó duy trì hành vi mở rộng giống A*, nhưng chỉ lưu trữ một đường đi duy nhất đang được khám phá cùng với 
    \textbf{giá trị $f$ tốt nhất} của các nhánh thay thế (alternative path cost).
    \item Khi đường đi hiện tại trở nên tệ hơn giới hạn $f_{limit}$, RBFS sẽ \textbf{quay lui (backtrack)} 
    và tiếp tục tìm kiếm ở nhánh có $f$ tốt nhất tiếp theo.
    \item Trực giác: RBFS hoạt động như A* “đệ quy”, nhưng thay vì lưu toàn bộ hàng đợi ưu tiên (frontier), 
    nó chỉ nhớ giá trị $f$ nhỏ nhất của các nhánh chưa khám phá sâu hơn — giúp giảm đáng kể bộ nhớ.
\end{itemize}

\textbf{Cấu trúc dữ liệu:} 
đệ quy trên cây tìm kiếm, biến \texttt{f\_limit}, ánh xạ \texttt{visited[node] = parent}, hàm heuristic \texttt{h(n)}.

\textbf{Cơ chế thực hiện:}
\begin{itemize}
    \item Giống A*, RBFS sử dụng hàm đánh giá:
    \[
        f(n) = g(n) + h(n)
    \]
    trong đó $g(n)$ là chi phí thực đã đi, còn $h(n)$ là ước lượng chi phí còn lại.
    \item Hàm RBFS(n, $f_{limit}$) hoạt động đệ quy:
    \begin{enumerate}
        \item Nếu $n$ là mục tiêu, trả về thành công.
        \item Sinh các nút con của $n$, tính $f$ cho mỗi nút:
        \[
            f(child) = \max(g(child) + h(child), f(n))
        \]
        để duy trì tính đơn điệu khi quay lui.
        \item Nếu không có con, trả về thất bại với giá trị $\infty$.
        \item Chọn nút con có $f$ nhỏ nhất làm “best child” và lưu nút có $f$ nhỏ thứ hai làm “alternative”.
        \item Nếu $f(\text{best}) > f_{limit}$, quay lui và trả về $f(\text{best})$.
        \item Ngược lại, gọi đệ quy:
        \[
            RBFS(\text{best}, \min(f_{limit}, f(\text{alternative})))
        \]
        để giới hạn tìm kiếm trong ngưỡng tốt nhất hiện tại.
    \end{enumerate}
    \item Khi quay lui, cập nhật giá trị $f$ tốt nhất của nhánh vừa duyệt để so sánh trong các lần quay lại.
\end{itemize}

\textbf{Ghi chú triển khai:}
\begin{itemize}
    \item RBFS là \textbf{thuật toán đệ quy theo chiều sâu}, sử dụng \textbf{bộ nhớ tuyến tính $O(bd)$}, 
    trong đó $b$ là hệ số phân nhánh, $d$ là độ sâu của lời giải.
    \item Mỗi lần quay lui, RBFS chỉ cần nhớ \textbf{giá trị $f$ tốt nhất} trong các nhánh chưa được khám phá sâu hơn.
    \item Do không lưu toàn bộ hàng đợi như A*, RBFS có thể \textbf{mở lại (re-expand)} các nút cũ khi quay lui, 
    khiến thời gian chạy có thể dài hơn.
\end{itemize}

\textbf{Đánh giá hiệu suất:}
\begin{itemize}
    \item \textbf{Tính đầy đủ (Completeness):} Có, nếu hệ số phân nhánh $b$ là hữu hạn và mọi chi phí cung đều dương.
    \item \textbf{Tính tối ưu (Optimality):} Có, nếu hàm heuristic $h(n)$ là \textbf{admissible} (khả chấp) — tương tự A*.
    \item \textbf{Độ phức tạp thời gian (Time Complexity):} Phụ thuộc vào độ chính xác của heuristic và tần suất quay lui.  
    Trường hợp xấu nhất, RBFS có thể có độ phức tạp gần $O(b^d)$ như A*.
    \item \textbf{Độ phức tạp không gian (Space Complexity):} $O(bd)$ — chỉ cần lưu đường đi hiện tại và các giá trị $f$ thay thế, 
    nhỏ hơn rất nhiều so với A* ($O(b^d)$).
    \item \textbf{Nhận xét:} RBFS đạt sự cân bằng tốt giữa tính tối ưu của A* và khả năng tiết kiệm bộ nhớ của DFS, 
    phù hợp cho các bài toán tìm kiếm sâu trong không gian lớn.
\end{itemize}